%! Setting Up a Test for the Difference of Two Population Proportions — Unit 6, Lesson 6.10 — Student Packet Cover Sheet
\documentclass[10pt]{article}
\usepackage{apstats-article}
\usepackage{apstats-boxes}
\usepackage{ltablex}
\keepXColumns

\begin{document}

% Full-bleed navy banner
\begin{tikzpicture}[remember picture, overlay]
  \fill[navy] (current page.north west)
    rectangle ([yshift=-0.9in]current page.north east);
\end{tikzpicture}
\vspace{-0.8in}

\begin{center}
  {\color{white}\textbf{\LARGE AP Statistics}}\\[4pt]
  {\color{skymid}\large\textbf{Unit 6}}\\[6pt]
  {\color{white}\textbf{\large Lesson 6.10 \quad Setting Up a Test for the Difference of Two Population Proportions}}
\end{center}

\vspace{0.22in}
\namedateperiod
\vspace{0.18in}

\begin{learningtargetbox}
\small
\begin{enumerate}[leftmargin=1.5em, itemsep=5pt]
  \item I can state $H_0$ and $H_a$ for a two-sample $z$-test about the difference of two population proportions. \textit{(VAR-6.H--I, Skill 1.E)}
  \item I can compute the pooled proportion $\hat p_c$ and use it to verify conditions (Random, Large Counts, 10\%) for a two-sample test. \textit{(VAR-6.J, Skill 4.C)}
\end{enumerate}
\end{learningtargetbox}

\vspace{0.14in}

\begin{tocbox}
\small
\renewcommand{\arraystretch}{1.5}
\begin{tabularx}{\linewidth}{@{} c l X r @{}}
\rowcolor{sky}
\textbf{\#} & \textbf{Component} & \textbf{Description} & \textbf{Score} \\
\midrule
1 & Warm-Up        & Review one-sample test setup; introduce pooling concept & \blank{1.2cm} \\
2 & Guided Notes   & Hypotheses, pooled $\hat p_c$, conditions, worked example & \blank{1.2cm} \\
3 & Group Activity & State hypotheses and verify conditions for two scenarios & \blank{1.2cm} \\
4 & Exit Ticket    & Full State + Plan for a two-sample test & \blank{1.2cm} \\
5 & Homework       & Two complete State + Plan problems & \blank{1.2cm} \\
\midrule
  & \multicolumn{2}{r}{\textbf{Total}} & \blank{1.2cm} \\
\end{tabularx}
\end{tocbox}

\end{document}
